% This is `mistex.tex' for use with MLTeX v2.2 (as of December 17, 1995).
%
% Copyright (C) 2000 by Thierry Laronde; all rights are reserved.
%
% MLTeX is copyright (C) 1990-92 by Michael J. Ferguson; all rights are reserved.
% MLTeX Version 2.2 is copyright (C) 1995 by B. Raichle; all rights are reserved.
% The TeX program is copyright (C) 1982 by D. E. Knuth.
% TeX is a trademark of the American Mathematical Society.
%
%
% Usage:
%
%  1. Run iniTeX with the MLTeX extensions on this file.
%     Do not try to use plain-TeX or LaTeX.
%
% Input of plain.tex
\input plain.tex
% The hyphen.tex loaded with plain.tex defines the hyphenation patterns and
% exceptions for the language0, in this case the english.
\def\english{\language=0}
% define correct substitutions
\charsubdef "C0 = "12 `A		% A grave
\charsubdef "C1 = "13 `A		% A accute
\charsubdef "C2 = "5E `A		% A circumflex
\charsubdef "C3 = "7E `A		% A tilde
\charsubdef "C4 = "7F `A		% A dieresis
\charsubdef "C5 = "17 `A		% A ring
\charsubdef "C7 = "18 `C		% C cedilla
\charsubdef "C8 = "12 `E		% E grave
\charsubdef "C9 = "13 `E		% E accute
\charsubdef "CA = "5E `E		% E circumflex
\charsubdef "CB = "7F `E		% E dieresis
\charsubdef "CC = "12 `I		% I grave
\charsubdef "CD = "13 `I		% I accute
\charsubdef "CE = "5E `I		% I circumflex
\charsubdef "CF = "7F `I		% I dieresis
\charsubdef "D1 = "7E `N		% N tilde
\charsubdef "D2 = "12 `O		% O grave
\charsubdef "D3 = "13 `O		% O accute
\charsubdef "D4 = "5E `O		% O circumflex
\charsubdef "D5 = "7E `O		% O tilde
\charsubdef "D6 = "7F `O		% O dieresis
\charsubdef "D9 = "12 `U		% U grave
\charsubdef "DA = "13 `U		% U accute
\charsubdef "DB = "5E `U		% U circumflex
\charsubdef "DC = "7F `U		% U dieresis
\charsubdef "DD = "13 `Y		% dagger or Y acute Capital
\charsubdef "E0 = "12 `a		% a grave
\charsubdef "E1 = "13 `a		% a accute
\charsubdef "E2 = "5E `a		% a circumflex
\charsubdef "E3 = "7E `a		% a tilde
\charsubdef "E4 = "7F `a		% a dieresis
\charsubdef "E5 = "17 `a		% a ring
\charsubdef "E7 = "18 `c		% c cedilla
\charsubdef "E8 = "12 `e		% e grave
\charsubdef "E9 = "13 `e		% e accute
\charsubdef "EA = "5E `e		% e circumflex
\charsubdef "EB = "7F `e		% e dieresis
\charsubdef "EC = "12 16		% i grave
\charsubdef "ED = "13 16		% i accute
\charsubdef "EE = "5E 16		% i circumflex
\charsubdef "EF = "7F 16		% i dieresis
\charsubdef "F1 = "7E `n		% n tilde
\charsubdef "F2 = "12 `o		% o grave
\charsubdef "F3 = "13 `o		% o accute
\charsubdef "F4 = "5E `o		% o circumflex
\charsubdef "F5 = "7E `o		% o tilde
\charsubdef "F6 = "7F `o		% o dieresis
\charsubdef "F9 = "12 `u		% u grave
\charsubdef "FA = "13 `u		% u accute
\charsubdef "FB = "5E `u		% u circumflex
\charsubdef "FC = "7F `u		% u dieresis
\charsubdef "FD = "13 `y		% double dagger or y acute small
\charsubdef "FF = "7F `y		% y dieresis
% We define now the 8bits characters so that they can be used for hyphenation.
% This implies defining the \lccode (lowercase character associated : the
% \patterns only uses lowercase comparison). 
% We define the \uccode as well (uppercase character associated).
\lccode"C0="E0 \uccode"C0="C0	% A grave
\lccode"C1="E1 \uccode"C1="C1	% A accute
\lccode"C2="E2 \uccode"C2="C2	% A circumflex
\lccode"C3="E3 \uccode"C3="C3	% A tilde
\lccode"C4="E4 \uccode"C4="C4	% A dieresis
\lccode"C5="E5 \uccode"C5="C5	% A ring
\lccode"C7="E7 \uccode"C7="C7	% C cedilla
\lccode"C8="E8 \uccode"C8="C8	% E grave
\lccode"C9="E9 \uccode"C9="C9	% E accute
\lccode"CA="EA \uccode"CA="CA	% E circumflex
\lccode"CB="EB \uccode"CB="CB	% E dieresis
\lccode"CC="EC \uccode"CC="CC	% I grave
\lccode"CD="ED \uccode"CD="CD	% I accute
\lccode"CE="EE \uccode"CE="CE	% I circumflex
\lccode"CF="EF \uccode"CF="CF	% I dieresis
\lccode"D1="F1 \uccode"D1="D1	% N tilde
\lccode"D2="F2 \uccode"D2="D2	% O grave
\lccode"D3="F3 \uccode"D3="D3	% O accute
\lccode"D4="F4 \uccode"D4="D4	% O circumflex
\lccode"D5="F5 \uccode"D5="D5	% O tilde
\lccode"D6="F6 \uccode"D6="D6	% O dieresis
\lccode"D9="F9 \uccode"D9="D9	% U grave
\lccode"DA="FA \uccode"DA="DA	% U accute
\lccode"DB="FB \uccode"DB="DB	% U circumflex
\lccode"DC="FC \uccode"DC="DC	% U dieresis
\lccode"DD="FD \uccode"DD="DD	% dagger or Y acute Capital
\lccode"E0="E0 \uccode"E0="C0	% a grave
\lccode"E1="E1 \uccode"E1="C1	% a accute
\lccode"E2="E2 \uccode"E2="C2	% a circumflex
\lccode"E3="E3 \uccode"E3="C3	% a tilde
\lccode"E4="E4 \uccode"E4="C4	% a dieresis
\lccode"E5="E5 \uccode"E5="C5	% a ring
\lccode"E7="E7 \uccode"E7="C7	% c cedilla
\lccode"E8="E8 \uccode"E8="C8	% e grave
\lccode"E9="E9 \uccode"E9="C9	% e accute
\lccode"EA="EA \uccode"EA="CA	% e circumflex
\lccode"EB="EB \uccode"EB="CB	% e dieresis
\lccode"EC="EC \uccode"EC="CC	% i grave
\lccode"ED="ED \uccode"ED="CD	% i accute
\lccode"EE="EE \uccode"EE="CE	% i circumflex
\lccode"EF="EF \uccode"EF="CF	% i dieresis
\lccode"F1="F1 \uccode"F1="D1	% n tilde
\lccode"F2="F2 \uccode"F2="D2	% o grave
\lccode"F3="F3 \uccode"F3="D3	% o accute
\lccode"F4="F4 \uccode"F4="D4	% o circumflex
\lccode"F5="F5 \uccode"F5="D5	% o tilde
\lccode"F6="F6 \uccode"F6="D6	% o dieresis
\lccode"F9="F9 \uccode"F9="D9	% u grave
\lccode"FA="FA \uccode"FA="DA	% u accute
\lccode"FB="FB \uccode"FB="DB	% u circumflex
\lccode"FC="FC \uccode"FC="DC	% u dieresis
\lccode"FD="FD \uccode"FD="DD	% double dagger or y acute small
\lccode"FF="FF \uccode"FF="DF	% y dieresis

% Additionnal characters : french guillemets
\font\threemmi=cmmi3
\catcode171=13	\def^^ab{{\threemmi \hbox{\raise1pt\hbox{<\kern-.4em <\kern-2pt}}}}         % 171 '253, "ab, 
\catcode187=13 	\def^^bb{{\threemmi \hbox{\raise1pt\hbox{\kern-2pt>\kern-.4em >\kern.25em}}}}       % 187, '273, "bb 
% \catcode171=13	\def^^ab{{\hbox{\raise.5ex\hbox{$<$\kern-.2em$<$\kern-.2em}}}}         % 171 '253, "ab, 
% \catcode187=13 	\def^^bb{{\hbox{\raise.5ex\hbox{\kern-.2em$>$\kern-.2em$>$\kern.2em}}}}       % 187, '273, "bb 
% Definition de \MisTeX
\def\MisTeX{M\raise.2ex\hbox{i}\lower.2ex\hbox{s}\kern-.1em\TeX}
% Hyphenation patterns and exceptions for french
\language1
\def\french{\language1\message{Setting language to 1 (french)}}
% Hyphenation patterns --- GUTenberg and others file
\input frhyph.tex
%% End of file `mistex.tex'.
